% GENERATED FILE. Edit canonical lessons, then run npm run generate:books.
% canonical-source-hash: fnv1a64:4a31ddcd9a961ef8
% canonical-lessons: RU-C467-moroz, RU-C467-gololyod, RU-C467-metel, RU-C467-myorznut, RU-C467-poskolznutsya

\chapter{Frost and ice}
\label{ch:ru-frost-and-ice}
\hlchaptermodality{467}

\begin{chapteropening}
\textbf{By the end of this chapter:} \emph{I can talk about the cold: frost, black ice, a blizzard, feeling the cold and slipping.}

Complete the last lesson of chapter 467: I can talk about the cold: frost, black ice, a blizzard, feeling the cold and slipping.
\end{chapteropening}

\section[\texorpdfstring{\ru{мороз} (moróz)}{мороз (moróz)}]{\ru{мороз} (moróz) — frost, freezing cold}

\label{lesson:RU-C467-moroz}



\begin{quote}
Before the new one: say the Russian for to rumble, then the Russian for to flash.
\end{quote}

\subsection*{You'll want to know: \ru{мороз}}
\textbf{\ru{мороз}} (\emph{moróz}) — \textquotedblleft{}frost, freezing cold\textquotedblright{}.

\begin{quote}
\ru{Утром} \ru{был} \ru{мороз}, \ru{и} \ru{окна} \ru{были} \ru{белые}.
\end{quote}

\begin{quote}
\emph{Útrom byl moróz, i ókna býli bélyye.}
\end{quote}

\begin{quote}
In the morning there was a frost, and the windows were white.
\end{quote}

Masculine. \emph{\ru{Минус} \ru{двадцать} — \ru{сильный} \ru{мороз}} — minus twenty is a hard frost.

\begin{cousinweb}
\emph{\ru{мороз}} comes from Old East Slavic \emph{morozŭ} (EtymDB), with the native \emph{-\ru{оро}-} shape; the same root gives \emph{\ru{мёрзнуть}} 'to freeze'.
\end{cousinweb}

\begin{grammarlens}[title={What you've built}]
One more word for telling, explaining and describing at B1.
\end{grammarlens}

\subsection*{Your turn}
\begin{itemize}
\raggedright
  \item \emph{Say it:} \emph{moróz}
  \item \emph{Say it:} \emph{moróz}, once more
  \item \emph{Say it:} a sentence of your own with \emph{moróz}
  \item \emph{Recall:} say the Russian for to rumble, then the Russian for to flash, then say \emph{moróz} again
\end{itemize}

\begin{tcolorbox}[breakable,colback=teal!4,colframe=teal!35!black,title={Before you move on}]
What does \ru{мороз} mean? (\textquotedblleft{}Frost, freezing cold\textquotedblright{}.) Say it once more.
\end{tcolorbox}
\section[\texorpdfstring{\ru{гололёд} (gololyód)}{гололёд (gololyód)}]{\ru{гололёд} (gololyód) — black ice}

\label{lesson:RU-C467-gololyod}



\begin{quote}
Before the new one: say the Russian for to flash, then the Russian for frost.
\end{quote}

\subsection*{You'll want to know: \ru{гололёд}}
\textbf{\ru{гололёд}} (\emph{gololyód}) — \textquotedblleft{}black ice\textquotedblright{}.

\begin{quote}
\ru{Из}-\ru{за} \ru{гололёда} \ru{машины} \ru{едут} \ru{медленно}.
\end{quote}

\begin{quote}
\emph{Iz-za gololyóda mashíny yédut médlenno.}
\end{quote}

\begin{quote}
Because of the black ice the cars are going slowly.
\end{quote}

Masculine. Built from \emph{\ru{голый}} — bare — and \emph{\ru{лёд}} — ice.

\begin{cousinweb}
\emph{\ru{гололёд}} is built from \emph{\ru{голый}} 'bare' and \emph{\ru{лёд}} 'ice': bare ice on the road. \emph{\ru{Лёд}} is inherited from Proto-Slavic \emph{ledŭ} (EtymDB).
\end{cousinweb}

\begin{grammarlens}[title={What you've built}]
One more word for telling, explaining and describing at B1.
\end{grammarlens}

\subsection*{Your turn}
\begin{itemize}
\raggedright
  \item \emph{Say it:} \emph{gololyód}
  \item \emph{Say it:} \emph{gololyód}, once more
  \item \emph{Say it:} a sentence of your own with \emph{gololyód}
  \item \emph{Recall:} say the Russian for to flash, then the Russian for frost, then say \emph{gololyód} again
\end{itemize}

\begin{tcolorbox}[breakable,colback=teal!4,colframe=teal!35!black,title={Before you move on}]
What does \ru{гололёд} mean? (\textquotedblleft{}Black ice\textquotedblright{}.) Say it once more.
\end{tcolorbox}
\section[\texorpdfstring{\ru{метель} (metél)}{метель (metél)}]{\ru{метель} (metél) — a blizzard, a snowstorm}

\label{lesson:RU-C467-metel}



\begin{quote}
Before the new one: say the Russian for frost, then the Russian for black ice.
\end{quote}

\subsection*{You'll want to know: \ru{метель}}
\textbf{\ru{метель}} (\emph{metél}) — \textquotedblleft{}a blizzard, a snowstorm\textquotedblright{}.

\begin{quote}
\ru{Из}-\ru{за} \ru{метели} \ru{мы} \ru{два} \ru{дня} \ru{сидели} \ru{дома}.
\end{quote}

\begin{quote}
\emph{Iz-za metéli my dva dnya sidéli dóma.}
\end{quote}

\begin{quote}
Because of the blizzard we stayed at home for two days.
\end{quote}

Feminine. Built on \emph{\ru{мести}} — to sweep.

\begin{cousinweb}
\emph{\ru{метель}} is built on \emph{\ru{мести}} 'to sweep': wind that sweeps the snow along. EtymDB lists \emph{\ru{метель}} without a deeper source.
\end{cousinweb}

\begin{grammarlens}[title={What you've built}]
One more word for telling, explaining and describing at B1.
\end{grammarlens}

\subsection*{Your turn}
\begin{itemize}
\raggedright
  \item \emph{Say it:} \emph{metél}
  \item \emph{Say it:} \emph{metél}, once more
  \item \emph{Say it:} a sentence of your own with \emph{metél}
  \item \emph{Recall:} say the Russian for frost, then the Russian for black ice, then say \emph{metél} again
\end{itemize}

\begin{tcolorbox}[breakable,colback=teal!4,colframe=teal!35!black,title={Before you move on}]
What does \ru{метель} mean? (\textquotedblleft{}A blizzard, a snowstorm\textquotedblright{}.) Say it once more.
\end{tcolorbox}
\section[\texorpdfstring{\ru{мёрзнуть} (myórznut)}{мёрзнуть (myórznut)}]{\ru{мёрзнуть} (myórznut) — to feel the cold, to freeze}

\label{lesson:RU-C467-myorznut}



\begin{quote}
Before the new one: say the Russian for black ice, then the Russian for a blizzard.
\end{quote}

\subsection*{You'll want to know: \ru{мёрзнуть}}
\textbf{\ru{мёрзнуть}} (\emph{myórznut}) — \textquotedblleft{}to feel the cold, to freeze\textquotedblright{}.

\begin{quote}
\ru{Без} \ru{шапки} \ru{уши} \ru{быстро} \ru{мёрзнут}.
\end{quote}

\begin{quote}
\emph{Bez shápki úshi býstro myórznut.}
\end{quote}

\begin{quote}
Without a hat your ears soon get cold.
\end{quote}

Imperfective; perfective \emph{\ru{замёрзнуть}}. Past \emph{\ru{мёрз}, \ru{мёрзла}}.

\begin{cousinweb}
\emph{\ru{мёрзнуть}} is inherited from Proto-Slavic \emph{mĭrznǫti} (EtymDB), on the same root as \emph{\ru{мороз}} 'frost'.
\end{cousinweb}

\begin{grammarlens}[title={What you've built}]
One more word for telling, explaining and describing at B1.
\end{grammarlens}

\subsection*{Your turn}
\begin{itemize}
\raggedright
  \item \emph{Say it:} \emph{myórznut}
  \item \emph{Say it:} \emph{myórznut}, once more
  \item \emph{Say it:} a sentence of your own with \emph{myórznut}
  \item \emph{Recall:} say the Russian for black ice, then the Russian for a blizzard, then say \emph{myórznut} again
\end{itemize}

\begin{tcolorbox}[breakable,colback=teal!4,colframe=teal!35!black,title={Before you move on}]
What does \ru{мёрзнуть} mean? (\textquotedblleft{}To feel the cold, to freeze\textquotedblright{}.) Say it once more.
\end{tcolorbox}
\section[\texorpdfstring{\ru{поскользнуться} (poskolznútsya)}{поскользнуться (poskolznútsya)}]{\ru{поскользнуться} (poskolznútsya) — to slip}

\label{lesson:RU-C467-poskolznutsya}



\begin{quote}
Before the new one: say the Russian for a blizzard, then the Russian for to feel the cold.
\end{quote}

\subsection*{You'll want to know: \ru{поскользнуться}}
\textbf{\ru{поскользнуться}} (\emph{poskolznútsya}) — \textquotedblleft{}to slip\textquotedblright{}.

\begin{quote}
\ru{На} \ru{гололёде} \ru{бабушка} \ru{поскользнулась} \ru{и} \ru{упала}.
\end{quote}

\begin{quote}
\emph{Na gololyóde bábushka poskolznúlas i upála.}
\end{quote}

\begin{quote}
Grandma slipped on the black ice and fell.
\end{quote}

Perfective; imperfective \emph{\ru{поскальзываться}}.

\begin{cousinweb}
\emph{\ru{поскользнуться}} is built from \emph{\ru{по}-}, \emph{\ru{скользнуть}} 'to slide once', from \emph{\ru{скользить}} 'to slide', and \emph{-\ru{ся}}. EtymDB lists the verb without a deeper source.
\end{cousinweb}

\begin{grammarlens}[title={What you've built}]
One more word for telling, explaining and describing at B1.
\end{grammarlens}

\subsection*{Your turn}
\begin{itemize}
\raggedright
  \item \emph{Say it:} \emph{poskolznútsya}
  \item \emph{Say it:} \emph{poskolznútsya}, once more
  \item \emph{Say it:} a sentence of your own with \emph{poskolznútsya}
  \item \emph{Recall:} say the Russian for a blizzard, then the Russian for to feel the cold, then say \emph{poskolznútsya} again
\end{itemize}

\begin{tcolorbox}[breakable,colback=teal!4,colframe=teal!35!black,title={Before you move on}]
What does \ru{поскользнуться} mean? (\textquotedblleft{}To slip\textquotedblright{}.) Say it once more.
\end{tcolorbox}
